% 2026-11-23 第 6 部：Post 429--443 頁（Post 第 I--VII 部）。註號與頁碼已以 PDF 頁圖核對（PDF 第 4、8、13、14 頁）；Post 註 6 引 Michelman「at 430」，所引定義實在 Michelman 400 頁（見 00_note），正文逐字引 Post 之引文並註 Post 頁。
% Rawls、Habermas、MacIntyre、Schauer、Fuller、Rousseau、Dewey、Frankfurter 及 Post 之《Constitutional Domains》等皆未取得藏本，一律轉引 Post 之註。
% Michelman 之引文頁（400、415、419、420、424、425--426、427）依 00_note 已回核 DASH 本。
\chapter{Post 之回應}
\label{ch:60_post}

\begin{qaquestion}
Post\ix{persons}{post@Post, Robert} 回應 Michelman\ix{persons}{michelman@Michelman, Frank I.} 的文章（429--443 頁），分七節。
\begin{enumerate}
\tightlist
\item Post 以什麼設例說明憲政主義\ix{terms}{03@\zhlidxcat{憲政主義與憲法之效力基礎}!憲政主義}與民主\ix{terms}{01@\zhlidxcat{民主理論}!民主}的緊張並非在一切情形成立？
\item Post 如何以 Rousseau\ix{persons}{rousseau@Rousseau, Jean-Jacques} 的普遍意志\ix{terms}{01@\zhlidxcat{民主理論}!普遍意志／個別意志}、Fuller\ix{persons}{fuller@Fuller, Lon L.} 的對抗制\ix{terms}{07@\zhlidxcat{法院之組織與功能}!對抗制}類比與社會分歧，反駁 Michelman 模型？
\item Post 所分的人民主權\ix{terms}{01@\zhlidxcat{民主理論}!人民主權}與民主各指什麼？法院的憲法解釋\ix{terms}{08@\zhlidxcat{憲法解釋與司法自制}!解釋}既不是民主自治\ix{terms}{01@\zhlidxcat{民主理論}!自治}，正當性\ix{terms}{02@\zhlidxcat{正當性與合理歧見}!正當性}來自何處？
\item Post 如何看待 Brennan\ix{persons}{brennan@Brennan, William J., Jr.} 的成就？
\end{enumerate}
\end{qaquestion}

\section{憲法理論之任務與廢除第 26 修正案之設例}
\label{sec:60_post-1}

Post 在 429 頁開頭主張，憲法理論的工作是揭示並釐清憲法裁判實踐內在的原則，因此不能站在 Rawls\ix{persons}{rawls@Rawls, John} 的「veil of ignorance」〔無知之幕〕之後，也不能置身 Habermas\ix{persons}{habermas@Habermas, Jürgen} 的「ideal speech situation」〔理想言說情境〕；Brennan 這類受敬重的大法官之成就，屬於待解釋的資料，Michelman 以 Brennan 的作品為探究的中心，在這個意義上是正確的。\footcite[429 \& nn.1--2]{post1998}Post 接受 Michelman 對兩個概念的界定：憲政主義是控制何種法律得以制定、經由何種程序制定的「law of lawmaking」〔造法之法\ix{terms}{03@\zhlidxcat{憲政主義與憲法之效力基礎}!造法之法}〕，民主是「[p]opular political self-government」〔人民的政治自治〕，即一國人民自行決定規範其政治結合之法律內容。\footcite[430 \& nn.5--6]{post1998}

問題出在 Michelman 稱民主在一個決定空間內受到「necessity」〔必然性〕的排除，這個空間是決定一國最基本之法、造法之法的內容；\footcite[430 \& n.7]{post1998}Michelman 原文見 420 頁。\footcite[420]{michelman1998}Post 認為憲法裁判實踐中有典型事例，使憲法意義服從於人民意志\ix{terms}{01@\zhlidxcat{民主理論}!人民意志}的專斷指示。他設想美國聯邦憲法經兩院各三分之二提案、四分之三各州議會批准，修正案廢除第 26 修正案\ix{statutes}{02@美國!01@聯邦憲法!126@增修條文第 26 條}，明文規定未滿 21 歲者不得投票選舉任何聯邦官員，一名 18 歲公民於是向聯邦法院起訴，請求投票權。\footcite[430--431 \& n.8]{post1998}Post 稱之為 paradigmatically easy case〔典型的易案\ix{terms}{08@\zhlidxcat{憲法解釋與司法自制}!易案}〕，法官若判原告勝訴，可謂未能理解憲法裁判的實踐，不問她個人的信念；她縱使認為這項修正案悖離民主，仍須以其反映依第 5 條\ix{statutes}{02@美國!01@聯邦憲法!005@第 5 條}程序所測得的人民意志為由駁回。結論是「popular sovereignty manifestly trumps constitutionalism」〔人民主權明顯勝過憲政主義〕。\footcite[431 \& n.9]{post1998}Post 於註 9 引 Schauer\ix{persons}{schauer@Schauer, Frederick} 的 Easy Cases 論易案。

Post 隨即自承「loaded the dice」〔在骰子上動手腳〕：設例選用的修正案既新又文義明確，且該修正案雖是對民主原則的一種解釋，並未與基本而普遍接受的民主原則直接衝突。他在註 10 另設一例，修正案授權總統終身任職，那會暴露他在原文第 V 部討論的民主與人民主權之緊張。\footcite[431 n.10]{post1998}Michelman 與多數憲法理論家所處理的案件則不同，涉及「cruel and unusual punishments」〔殘酷而異常之刑罰〕、「the equal protection of the laws」〔法律之平等保護〕、「due process of law」〔正當法律程序〕、「the freedom of speech」〔言論自由〕這類古老而文字概括的條款。Post 同意 Michelman 的結論：歷史使人難以把這些條款的解釋與適用視為帶有人民意志的權威，縱使條款曾經由人民制定；只要這些條款仍能作為撤銷一般立法所表達之人民意志的理由，Michelman 所指出的憲政主義與民主之緊張就無從消解。\footcite[431 \& n.11]{post1998}

\section{Michelman 模型與 Rousseau 之普遍意志}
\label{sec:60_post-2}

Michelman 在文末提出一項「surmise」〔推測〕，說明 Brennan 可能如何解決這項緊張：出於對法律的「respect」〔尊重〕而自願遵守者，若在反思後認為法律公正，即處於「in a state of self-government, in as full a sense as we can reasonably hope to find in politics」〔處於我們在政治中所能合理期待找到的最完整意義上的自治狀態〕。\footcite[431--432 \& n.12]{post1998}Post 說他不清楚 Michelman 本人是否接受這項立場，有文本證據顯示 Michelman 並不完全滿意，Post 自己的主張則是「definitively repudiated」〔應予明確駁斥〕。\footcite[432]{post1998}

這項立場預設「procedure-independent standards of rightness」〔程序獨立之正確標準\ix{terms}{02@\zhlidxcat{正當性與合理歧見}!程序獨立之正確標準}〕，對憲法判斷的尊重因而不能只取決於測量人民意志\ix{terms}{01@\zhlidxcat{民主理論}!人民意志}的程序性標準，例如使 Post 所設修正案生效的第 5 條。Post 在 432 頁引 Michelman 425--426 頁所述令人尊重的根本法\ix{terms}{03@\zhlidxcat{憲政主義與憲法之效力基礎}!根本法}須具備的兩項屬性：其一，須使人看到有人持續盡「maximum feasible effort」〔最大可行的努力〕把根本法與其主要解釋弄對，努力可包括設立專司解釋的法庭；其二，須有安排使有權決定根本法的人暴露於社會上「the full blast of sundry opinions and interest-articulations」〔各種意見與利益表達的全面衝擊〕，其中對自己意見的登記所以重要，在於其作為辯論貢獻的「presumptive epistemic value」〔推定的認識價值\ix{terms}{02@\zhlidxcat{正當性與合理歧見}!認識價值}〕，與影響決定的槓桿無關。\footcite[432]{post1998}\footcite[425--426]{michelman1998}Post 僅為行文方便，稱這兩項條件為「Michelman's model」。

這個模型以大法官的判斷被普遍接受為「just, or fair」或「right」為前提。Post 指出這個前提回到 Rousseau\ix{persons}{rousseau@Rousseau, Jean-Jacques} 以來集體自治的核心問題：在顯著而無從消除的歧異之下，如何說我們經由集體自決而自治，多數決的每一次制定為何不同時是對少數的壓迫。Rousseau 的回答是使每個公民的個別意志與全體的普遍意志\ix{terms}{01@\zhlidxcat{民主理論}!普遍意志／個別意志}融合，服從普遍意志者「obeys no one but himself, and remains as free as before」〔只服從自己，而且與從前一樣自由〕。\footcite[433 \& nn.15--16]{post1998}Post 認為，Michelman 假定大法官的判斷被接受為正確，便是假定與 Rousseau 同一種個別意志與普遍意志的融合：同意 Brennan 這樣的大法官正確解釋了根本法的人，其個別意志即與表現於 Brennan 判決的普遍意志相符，服從這些判決等於服從自己。\footcite[433 \& n.17]{post1998}

對 Rousseau 的指責是假定融合而無視現實，Post 認為 Michelman 模型同樣承受這項批評。他承認自己對 Brennan 的評價極高，註 18 引他自己論 Brennan 與 Warren 法院的文章；\footcite[433 n.18]{post1998}但 Brennan 的作品仍有爭議，大部分人口不接受其為「just」或「right」，任何大法官皆然，這是 Rawls 所稱現代社會「diversity of conflicting and irreconcilable---and what's more, reasonable---comprehensive doctrines」〔相互衝突、無可調和而且合理的整全性學說\ix{terms}{02@\zhlidxcat{正當性與合理歧見}!整全性學說}之多樣性〕的結果。\footcite[433 \& n.19]{post1998}對同意某項憲法解釋的個人，這個模型可以說明他如何認同該判斷，卻不能說明一個多元的國家如何整體認同。\footcite[433--434]{post1998}

\section{對抗制之類比與 Taft Court 時期之批評}
\label{sec:60_post-3}

Post 在 434 頁先敘述 Rousseau\ix{persons}{rousseau@Rousseau, Jean-Jacques} 的近代批評者如何改走另一條路：既然個別意志\ix{terms}{01@\zhlidxcat{民主理論}!普遍意志／個別意志}無法確定地等同於國家的各個制定，就把個人意志連到制定法律的體系，使立法取決於公眾意見\ix{terms}{01@\zhlidxcat{民主理論}!公眾意見}，並以人人得自由參與形成公眾意見的公共討論，為公民認同公眾意見的必要條件，Post 在《Constitutional Domains》184--187 頁有同樣的論述。\footcite[434 \& n.20]{post1998}Michelman 模型也暗示一套確保民主自決而無須對特定法律內容達成共識的體系，其兩項條件是專司解釋的法庭，以及使法庭成員暴露於社會各種意見之下的安排。\footcite[434 \& n.21]{post1998}前一種體系著重公民參與形成所受治之法，Michelman 模型明白排除這個觀念，Post 因此把問題界定為：缺少對個別憲法判斷結果之共識時，Michelman 所提的結構能否調和憲政主義與民主。\footcite[434]{post1998}

Post 對此的分析從 Fuller\ix{persons}{fuller@Fuller, Lon L.} 對對抗制\ix{terms}{07@\zhlidxcat{法院之組織與功能}!對抗制}的詮釋入手。對抗制同樣有專事解釋與適用法律的法庭，也設置辯護人，使法院下判斷前雙方的主張「be carefully considered and be given its full weight and value」〔得到仔細考慮並被給予完整的分量與價值〕，Fuller 引自 Talks on American Law 所收〈The Adversary System〉，Post 於註 22、23 引用。\footcite[434--435 \& nn.22--23]{post1998}Post 願意承認這些條件大致確立了對抗制所為判斷的正當性，Michelman 有時也像在談「the problem of political legitimacy」〔政治正當性\ix{terms}{02@\zhlidxcat{正當性與合理歧見}!政治正當性}的問題〕。\footcite[435 \& n.24]{post1998}但據 Post 所知，沒有人，Fuller 更不曾，主張這些條件足以把對抗制稱為民主自治的結構。訴訟當事人受法院管轄，卻無人說他們經由行使管轄而自治；若要主張，也只能援引當事人以公民身分參與了法院所適用之法律的形成，Michelman 模型所要排除的正是這項參與。\footcite[435]{post1998}

Post 由此推論，全國人民也不能說經由最高法院的憲法解釋而自治。1920 年代 Taft Court 大力恢復實質正當程序，不同意其憲法解釋者並不視之為自治的工具，而攻擊它是「Unreformed House of Lords」〔未經改革的上議院〕，見 1922 年 6 月 16 日 The N.Y. Call 第 1 版之 Supreme Court Nullifies People's Will 一文；是「absolute despotism」〔絕對專制〕，見 1924 年 8 月 31 日 N.Y. Times 第 1 版之 Justice Ford Backs La Follette Ticket 一文；並且「militates against every conception upon which democratic governments are supposed to be based」〔違背一切民主政府所應賴以建立的觀念〕，見 Meyer London 刊於 1923 年 3 月號 American Federationist 225 頁的文章。\footcite[435 \& nn.25--27]{post1998}Post 預期讀者覺得這些措詞誇大，因而請讀者設想一個當代法院，其大法官能幹而盡職，受到「the full blast of sundry opinions and interest-articulations in society」，卻一貫採取與讀者自己深刻對立的法學見解，亦即由九位 McReynolds\ix{persons}{mcreynolds@McReynolds, James C.} 組成的法院；他很懷疑讀者會把這樣的法院視為自治的工具，縱使 Michelman 的兩項結構條件都已滿足。\footcite[435]{post1998}

Post 的結論是 Michelman 模型取決於法院把主要解釋弄對；不相信法院正確解釋憲法時，那兩項結構條件本身無法產生民主自決的價值，「At most, these conditions can establish only political legitimacy」〔至多，這些條件只能確立政治正當性〕。\footcite[436]{post1998}

\section{公眾意見體系與 Michelman 模型下之社會分歧}
\label{sec:60_post-4}

Post 肯定 Michelman 提出「getting it right」〔把它弄對〕與民主自治之關係的問題是正確的。即便立法取決於公眾意見\ix{terms}{01@\zhlidxcat{民主理論}!公眾意見}，人人得自由參與形成公眾意見，一個持續被多數擊敗、忽視而疏離的群體，仍可能質疑這個國家是否為民主自治的合適載體，甚至考慮脫離、另建民主政體；自由開放的民主參與為民主自治所必要，單憑它並不足夠。\footcite[436]{post1998}

公眾意見體系中，民主正當性\ix{terms}{02@\zhlidxcat{正當性與合理歧見}!民主正當性}獨立於該體系的特定立法結果，因而與 Rawls 所稱「irreparable reasonable disagreement」〔無從化解之合理歧見\ix{terms}{02@\zhlidxcat{正當性與合理歧見}!無從化解之合理歧見}〕大致相容；\footcite[436 \& n.28]{post1998}持續而根本的分歧是這項獨立性的「limiting condition」〔極限條件〕，標示民主與多樣性可能變得無從調和的極端點。\footcite[436]{post1998}Michelman 模型下，公民只在同意法院憲法判斷的範圍內才算自治，「Consensus defines and exhausts the scope of democratic self-governance」〔共識界定並窮盡民主自治的範圍〕；分歧於是不只在極端情形，而在一切情形都否定民主正當性，模型沒有辦法跨越民主正當性與現代社會之無從化解之合理歧見\ix{terms}{02@\zhlidxcat{正當性與合理歧見}!合理歧見}之間的鴻溝。\footcite[436--437]{post1998}Post 把這點列為拒絕該模型作為當代國家民主自治之一般說明的有力理由。\footcite[437]{post1998}

民主自治的一般說明並非 Michelman 計畫的焦點。Michelman 在較窄的脈絡中發展模型，要調和民主與違憲審查\ix{terms}{04@\zhlidxcat{違憲審查權限之歸屬}!違憲審查}所體現的憲政主義，並採取 Post 認為準確的嚴苛假設：多數憲法解釋不能合理地解為反映公眾意見，一般的民主參與體系因此不能賦予某項憲法判斷以民主正當性。Post 接下來所評估的，是 Michelman 歸於 Brennan 的「surmise」在這個狹窄得多的脈絡中的表現。\footcite[437]{post1998}

\section{人民主權與民主之區分}
\label{sec:60_post-5}

Post 以美國的憲政主義必然包含違憲審查\ix{terms}{04@\zhlidxcat{違憲審查權限之歸屬}!違憲審查}為出發點：法院判斷是否撤銷原應反映人民意志\ix{terms}{01@\zhlidxcat{民主理論}!人民意志}的合法制定。第 1 修正案\ix{statutes}{02@美國!01@聯邦憲法!101@增修條文第 1 條}的判決常以該立法侵害「vital to the maintenance of our democratic institutions」〔對維持我國民主制度至關重要〕的表現自由為由撤銷，見 Schneider v. State\ix{cases}{02@美國!01@聯邦最高法院!1939Schneider v. State（1939）@Schneider v. State（1939）}, 308 U.S. 147, 161 (1939)；第 1 修正案並被稱為「the guardian of our democracy」〔我們民主的守護者〕，見 Brown v. Hartlage\ix{cases}{02@美國!01@聯邦最高法院!1982Brown v. Hartlage（1982）@Brown v. Hartlage（1982）}, 456 U.S. 45, 60 (1982)。\footcite[437 \& nn.29--30]{post1998}Post 由此推出，在美國的憲法傳統中，民主不能等同於人民意志的多數決表達。

憲法理論因而須區分民主與 Post 約定稱為 popular sovereignty〔人民主權〕的東西：國家服從於人民意志，而人民意志依多數決或第 5 條修憲機制等程序標準認定。Post 開頭所設的修正案例屬人民主權之例，未必是民主之例；人民若依法制定廢除投票權、授予聯邦官員終身世襲任期的修正案，便是人民主權之例，卻顯然違反民主。\footcite[437]{post1998}依 Post，人民主權依程序標準界定，民主依實質價值界定；他引 Michelman 的話：「Democracy is a demanding normative idea, an idea with content, however uncertain or disputable that content may be」〔民主是一個要求甚高的規範理念，一個有內容的理念，不論那內容多麼不確定或有爭議〕。\footcite[438 \& n.31]{post1998}\footcite[419]{michelman1998}Post 另文主張，在美國憲法傳統中，民主所以受重視，在於集體自治的價值，該價值有強大而可爭議的含義，包括一個不受檢查、永遠開放而預設自主公民參與的公共論述\ix{terms}{01@\zhlidxcat{民主理論}!公共論述}結構；\footcite[438 \& n.32]{post1998}註 33 說，Michelman 歸於他的程序性民主說，大致只表達了實質民主觀對公共論述規制的一項含義，並參 Dworkin\ix{persons}{dworkin@Dworkin, Ronald} 的 The Arduous Virtue of Fidelity，Post 引為較完整的說明。\footcite[438 n.33]{post1998}法院據第 1 修正案撤銷牴觸這些含義的多數決立法時，傳統上稱為民主與憲政主義之間的緊張，可更準確地分析為人民主權與憲法上之民主理想之間的衝突。\footcite[438]{post1998}

Post 同意 Michelman 的兩點：公眾意見\ix{terms}{01@\zhlidxcat{民主理論}!公眾意見}並不授權法院為評價無數無法預見的案件而作出的各項民主實質價值解釋，民主也無法「decide the question of what democracy means in the first place」〔決定民主一開始是什麼意思這個問題〕。\footcite[438 \& n.34]{post1998}由此，Michelman 迫使人們追問，適用民主實質價值的憲法判斷能否與民主所要求的自治相容。Post 認為，依 Michelman 模型，這類判斷只在且只在我們相信法院作成正確決定的範圍內是民主的，因為此時我們受自己所同意的法律拘束，這個答覆有說服力。但借 Michelman 的詞，這個立場傳達的「dark-hued message」〔色調暗沉的訊息〕是，在「irreparable reasonable disagreement」之下，適用民主價值的憲法裁判，對大片人口將始終不能與自治的實踐調和：被排除於所受治之法的形成之外、又不以其他方式認同該法的人，不能把該法歸為自己的民主創造。Post 引 Michelman 的話收尾：「To find the laws deserving of your respect is not yet to decide the laws」〔認為法律值得你尊重，尚不是決定法律〕。\footcite[438--439 \& n.35]{post1998}\footcite[427]{michelman1998}

\section{民主自治以外之正當性來源與社群}
\label{sec:60_post-6}

Post 在 439 頁指出，憲法對民主價值的解釋並不因此欠缺正當性，因為民主自決之外尚有許多正當性來源。民主的實踐嵌在非民主的社會組織形式之中，並預設它們持續繁榮；對集體自決價值的持續承諾，本身取決於社會化與文化再生產的過程，這些過程是非民主的，屬於 Post 稱為「community」〔社群\ix{terms}{01@\zhlidxcat{民主理論}!社群}〕的社會秩序\ix{terms}{01@\zhlidxcat{民主理論}!社會秩序}。\footcite[439 \& nn.36--37]{post1998}

Post 區分兩者的方向。民主面向未來，關乎一個民族決定自己命運的權力，因此要求社會結構持續保有重塑個人與集體認同的可能；社群則回溯過去，要求社會結構設法表述並強制那些被推定為早已構成個人與集體認同的社會規範。\footcite[439 \& n.38]{post1998}民主所設想的開放未來，仍然預設至少一項既存價值，即民主本身，因而民主依賴一個維持並滋養集體自決，使之構成集體與個人認同的社會結構。集體自決的價值於是是一項社群規範\ix{terms}{01@\zhlidxcat{民主理論}!社群規範}，如一切社群規範，經由社會化與制度化而再生產，本質上可爭議。\footcite[439 \& n.40]{post1998}我們必然在民主的意義上有歧見，因為對其意義採取立場，是在表達我們是怎樣一個民族的願景。\footcite[439]{post1998}

註 39 補充 Post 所稱 paradox of public discourse〔公共論述之悖論\ix{terms}{01@\zhlidxcat{民主理論}!公共論述之悖論}〕：在美國憲法傳統中，民主被詮釋為在公共論述\ix{terms}{01@\zhlidxcat{民主理論}!公共論述}之內中止執行那些維持我國文化對民主之實踐與持續承諾的社群規範。\footcite[439 n.39]{post1998}民主與人民主權的界線因此總有爭論餘地。Post 認為 Michelman 文章最突出的貢獻之一，是持續照亮民主自治理想同時要求人民主權與對其限制的悖論；這些限制最終所具的根據，不高於我們「effort to read the Constitution in the light of what John Dewey once called 'The Great Community' of the nation」〔依 John Dewey 所稱國家的「偉大社群」來解讀憲法的努力〕，Post 引自《Constitutional Domains》18 頁。\footcite[440 \& n.41]{post1998}

\section{國族精神與第 5 條之承認規則}
\label{sec:60_post-7}

Post 說這項分析不限於第 1 修正案與民主價值：憲法解釋的實踐本身並非建基於民主自決，民主自決也只是憲法解釋所促進的眾多價值之一。這些價值涵蓋聯邦制（與第 11 修正案文本相連）、實質平等（平等保護條款所蘊含）、公平（正當程序條款所體現）以及自由人道主義（通常位於第 8 修正案文本）。生動的憲法裁判多半要窺探「behind the words of the constitutional provisions」〔憲法條文用語之後〕，辨識「postulates which limit and control」〔設限與控制的前提〕，見 Principality of Monaco v. Mississippi\ix{cases}{02@美國!01@聯邦最高法院!1934Principality of Monaco v. Mississippi（1934）@Principality of Monaco v. Mississippi（1934）}, 292 U.S. 313, 322 (1934)；這些前提最終植根於「constitutional ideals」〔憲法理想\ix{terms}{03@\zhlidxcat{憲政主義與憲法之效力基礎}!憲法理想}〕，見 Planned Parenthood of Southeastern Pa. v. Casey\ix{cases}{02@美國!01@聯邦最高法院!1992Planned Parenthood of Southeastern Pa. v. Casey（1992）@Planned Parenthood of Southeastern Pa. v. Casey（1992）}, 505 U.S. 833, 868 (1992)（複數意見），即對國族社群\ix{terms}{01@\zhlidxcat{民主理論}!社群}歷史上之根本規範的解釋，合稱 national ethos〔國族精神\ix{terms}{03@\zhlidxcat{憲政主義與憲法之效力基礎}!國族精神}〕。\footcite[440 \& nn.42--44]{post1998}

註 44 舉出國族精神內各項理想可能彼此衝突的例子：民主的實質價值被詮釋為要求開放的公共論述\ix{terms}{01@\zhlidxcat{民主理論}!公共論述}結構，中止一切強加特定國族認同\ix{terms}{01@\zhlidxcat{民主理論}!國族認同}觀念的規制，包括表達相競憲法理想的規制；仇恨言論之爭常被表述為第 1 修正案與第 14 修正案的衝突，Post 據此認為傳統憲法理論所設想的民主與憲政主義之緊張，有時可更準確地看作憲法上民主理想與其他憲法理想的衝突。\footcite[440 n.44]{post1998}

人民主權與國族精神的關係，Post 這樣描述：司法對國族精神的闡述「flourishes in the interstices of popular sovereignty」〔在人民主權的縫隙中茁壯〕。憲法解釋隨時可被人民依第 5 條的程序性標準所測得的專斷意志修正，這些標準有如「rule of recognition」〔承認規則\ix{terms}{03@\zhlidxcat{憲政主義與憲法之效力基礎}!承認規則}〕，把民意轉成實證法，其強制的指令赤裸而不容置疑。但隨著法律的當時脈絡淡去、無法預見的適用累積並自成一格，違憲審查\ix{terms}{04@\zhlidxcat{違憲審查權限之歸屬}!違憲審查}歷來傾向把概括或模糊的憲法用語轉化為國族可爭議之理想的所在；我們趨向一種強調正當化實踐的法治觀，使憲法對人民主權變得「opaque」〔不透明〕。\footcite[440--441 \& n.45]{post1998}

註 45 針對同一個難處：「democracy」可以在人民主權之硬性實證論與民主自治之軟性規範之間滑動，「constitutionalism」也可以在第 5 條所定之法的硬性實證論與國族精神的軟性社會理想主義之間滑動，分析兩者的緊張時須先確定被對照的詞義。Post 以 2×2 表呈現，四格是民主的兩種意義與憲政主義的兩種意義之組合。\footcite[441 n.45]{post1998}

\begin{longtable}{@{}L{0.22\linewidth}L{0.36\linewidth}L{0.36\linewidth}@{}}
\caption{Post 註 45 之「Democracy versus Constitutionalism」〔民主對憲政主義〕四格表。}\\
\toprule
 & 憲政主義為國族精神（規範性）Constitutionalism as National Ethos (Normative) & 憲政主義為實證法 Constitutionalism as Positive Law \\
\midrule
\endfirsthead
\toprule
 & 憲政主義為國族精神（規範性）Constitutionalism as National Ethos (Normative) & 憲政主義為實證法 Constitutionalism as Positive Law \\
\midrule
\endhead
民主為規範理想 Democracy as Normative Ideal & 格 1：憲法上的民主理想與其他憲法理想之衝突（見註 44） & 格 2：實證憲法要求（例如第 5 條）與民主理想之關係 \\
民主為人民主權 Democracy as Popular Sovereignty & 格 3：憲法規範（例如第 1 修正案所含民主理想）與多數決立法之對立，屬例行的憲法裁判爭點 & 格 4：實證制度設計，例如多數決與特別多數決之利弊 \\
\bottomrule
\end{longtable}

\section{違憲審查對社群規範之闡述與法院之處境}
\label{sec:60_post-8}

Post 在 441 頁列出法院闡述憲法理想\ix{terms}{03@\zhlidxcat{憲政主義與憲法之效力基礎}!憲法理想}的手段：教育、說服、論證與啟發等熟悉的技術，其機構權威的相當聲望，以及傳統法律推理的強大推定；法院以這些資源鞏固既有的文化承諾，也促成新文化承諾的發展。這些場合中法院自己的憲法解釋，沒有一次實現民主自治的實踐，它們所展示的是法律如何表現並強制「ambient community norms」〔瀰漫的社群規範\ix{terms}{01@\zhlidxcat{民主理論}!社群規範}〕。\footcite[441 \& n.46]{post1998}在文化多元的情形中，法律作此功能時，對不分享法律所表述之文化願景者，有「hegemonic domination」〔霸權支配〕的可能。\footcite[441 \& n.47]{post1998}

在憲法裁判中，這種霸權有時是把憲法理想強加於持異議的地理區域，Post 舉 Brown v. Board of Education\ix{cases}{02@美國!01@聯邦最高法院!1954Brown v. Board of Education（1954）@Brown v. Board of Education（1954）}, 347 U.S. 483 (1954)，有時是把世代（或階級）的理想強加於敵對的人群，如新政時期的憲法危機。法院在反對聲中推行特定的國族精神\ix{terms}{03@\zhlidxcat{憲政主義與憲法之效力基礎}!國族精神}願景時，易受政治報復，形式有公民不服從、敵意的總統任命或憲法修正；通過這些危險之困難與必要，便是何以說最高法院大法官必須是「statesmen」〔政治家〕，Post 於註 49 引 Frankfurter\ix{persons}{frankfurter@Frankfurter, Felix} 的 The Court and Statesmanship 為例。\footcite[442 \& nn.48--49]{post1998}

\section{Brennan 作為司法政治家與 Post 之結論}
\label{sec:60_post-9}

Post 在 442 頁稱 Brennan 是偉大的 judicial statesman〔司法政治家\ix{terms}{08@\zhlidxcat{憲法解釋與司法自制}!司法政治家}〕：他深知自己權位的限度與義務，知道何時何處妥協，而最重要的是，他極有紀律而有力地闡述憲法理想\ix{terms}{03@\zhlidxcat{憲政主義與憲法之效力基礎}!憲法理想}，尤其是民主自治理想；他對第 1 修正案法學的影響「second to none」〔無人能出其右〕。Post 完全同意 Michelman 的評價：Brennan 認為法院被賦予權威與責任，為國家的政治體制解釋一項程序獨立的正確、正義與民主標準，其可貴在於他認為法院僅在行使權力時著眼於保護並擴大人人向法院與其他社會權威民主地表達意見的權利與機會，才享有這項授權。\footcite[442 \& n.50]{post1998}\footcite[426]{michelman1998}

Post 接著指出，Brennan 的成就不能被描述為「self-government through democracy」〔經由民主的自治〕，語出 Michelman 同頁。\footcite[442 \& n.51]{post1998}Brennan 對民主的憲法意義提出有說服力的說明，並成功地將其大部分轉成憲法；與 Brennan 觀點相同者有理由認同這部法，不同意者則否。Post 承認一項反諷：Brennan 的憲法判斷創造的各項自由，增進了所有人的民主自治，連不同意他的人也在內；這項反諷並不使這些判斷本身成為異議者之自治的範例。異議者因 Brennan 的工作而享有的自治，只是可用以爭取他們支持 Brennan 憲法願景的一項理由。\footcite[442--443]{post1998}

Post 以這項願景作結：它享有被（法律上）權威地表述的國族認同\ix{terms}{01@\zhlidxcat{民主理論}!國族認同}的地位，因此是國族持續的社群\ix{terms}{01@\zhlidxcat{民主理論}!社群}自我界定過程的一部分，我們全都牽涉其中，憲法理論也能在其中發揮微小的作用；他引自己〈Subsidized Speech〉192--193 頁對民主自治與社群自我界定之別的說明，並在文末讚揚 Brennan 法學的公正與深刻。\footcite[443 \& n.52]{post1998}Post 對 Michelman 模型的讀法，與 Michelman 原文不合之處，見第 \ref{ch:70_misreadings} 部。
